\subsection{Ograniczenia}

\begin{enumerate}
  \item \textbf{Liczba zbiorów.} Jednostką analizy jest zbiór danych, a~w~każdym zadaniu jest ich 12. Przy tej liczbie zbiorów i~korekcie Holma test Wilcoxona wykrywa tylko efekty obecne w~prawie wszystkich zbiorach (Tabela~\ref{tab:moc}). Wynik nieistotny oznacza więc brak wystarczających dowodów, a~nie brak efektu. Częściowo kompensują to analiza łączna 24 zbiorów i~test bayesowski, który osobno wskazuje przypadki praktycznej równoważności.
  \item \textbf{Korelacje między wariantami w~H2 odzwierciedlają plan eksperymentu.} 15 wariantów to zaplanowane punkty, a~nie losowa próba sposobów losowania. Miary takie jak $U$ i~$D_J$ są prawie wyznaczone przez parametry wariantu ($r$, $s$, $\gamma$). Korelacja miary z~jakością zespołu pokazuje więc, które własności sposobu losowania współwystępują z~jakością, a~nie że sama wartość miary ją powoduje. Korelacja cząstkowa usuwa tylko tę część związku, która przechodzi przez średnią jakość drzew, i~to w~zakresie zależności liniowej między rangami.
  \item \textbf{Koncentracja miar.} Dla ustalonego wariantu miary próbek zmieniają się między losowaniami o~wielkość rzędu $1/\sqrt{n}$ (Sekcja~\ref{sec:koncentracja}). Efekty w~obrębie jednego wariantu (H1, eksperyment dodatkowy) muszą więc być bardzo małe. Wynika to z~własności losowania, a~nie z~niedokładności pomiaru.
  \item \textbf{Ustalone hiperparametry.} Używamy drzew bez ograniczenia głębokości i~$M=100$. Przy drzewach przycinanych, które są stabilniejsze, wpływ losowania próbek byłby prawdopodobnie mniejszy, a~przy znacznie większym $M$ zmalałby udział losowego szumu w~wyniku zespołu.
  \item \textbf{Przygotowanie danych.} Usunięcie wierszy z~brakami w~Adult zakłada, że braki występują całkowicie losowo (MCAR). Duże zbiory ograniczyliśmy do 20\,000 obserwacji. Obie decyzje dotyczą wszystkich wariantów jednakowo, więc nie zaburzają ich porównania, ale zawężają populację, której dotyczą wnioski.
  \item \textbf{Zmienność między podziałami.} Odchylenie standardowe wyników z~powtarzanej walidacji krzyżowej zaniża rzeczywistą niepewność, bo części treningowe się pokrywają~\cite{nadeau2003}. Kryterium znaczenia praktycznego, czyli różnicę większą niż to odchylenie, traktujemy więc jako orientacyjne, a~wnioski opieramy na testach po zbiorach.
\end{enumerate}
